% !TeX program = lualatex
% Compile twice: lualatex 129.tex
% All references and prose are embedded; no BibTeX database or external inputs.
\documentclass[11pt,a4paper]{article}
\usepackage[margin=24mm,headheight=14pt]{geometry}
\usepackage{fontspec}
\setmainfont{Latin Modern Roman}
\setsansfont{Latin Modern Sans}
\setmonofont{Latin Modern Mono}
\usepackage{amsmath,amssymb,mathtools}
\usepackage{unicode-math}
\setmathfont{Latin Modern Math}
\usepackage{microtype}
\usepackage{enumitem,xurl,needspace}
\usepackage[dvipsnames]{xcolor}
\usepackage{hyperref}
\hypersetup{unicode=true,colorlinks=true,linkcolor=MidnightBlue,urlcolor=MidnightBlue,
 pdftitle={500 Mathematical Problem Statements},pdfauthor={Source-annotated compilation},bookmarksnumbered=true}
\usepackage{bookmark}
\usepackage{tocloft}
\setlength{\cftsecnumwidth}{3em}
\cftsetpnumwidth{2.5em}
\cftsetrmarg{4em}
\usepackage{titlesec}
\titleformat{\section}{\Large\bfseries\raggedright}{\thesection}{0.7em}{}
\usepackage{fancyhdr}
\pagestyle{fancy}\fancyhf{}
\fancyhead[L]{\small\sffamily Mathematical problem statements}
\fancyhead[R]{\small Source edition 22}
\fancyfoot[C]{\thepage}
\setlength{\parindent}{0pt}
\setlength{\parskip}{5pt plus 1pt}
\setlength{\emergencystretch}{3em}
\setcounter{tocdepth}{1}
\setcounter{secnumdepth}{1}
\setlist[itemize]{leftmargin=3em,itemsep=5pt,topsep=4pt,parsep=0pt}
\allowdisplaybreaks
\newcommand{\N}{\mathbb{N}}
\newcommand{\Z}{\mathbb{Z}}
\newcommand{\Q}{\mathbb{Q}}
\newcommand{\R}{\mathbb{R}}
\newcommand{\C}{\mathbb{C}}
\newcommand{\F}{\mathbb{F}}
\newcommand{\A}{\mathbb{A}}
\newcommand{\PP}{\mathbb P}
\newcommand{\E}{\mathbb{E}}
\newcommand{\eps}{\varepsilon}
\newcommand{\dd}{\,\mathrm d}
\newcommand{\sref}[2]{\hyperlink{source:#1:#2}{[S#2]}}
\newcommand{\eref}[2]{\hyperlink{extra:#1:#2}{[E#2]}}
% Turn the next switch off for a shorter reading copy. The quotations preserve
% every exactTarget field, including qualifications omitted from a short summary.
\newif\ifshowcatalogue
\showcataloguetrue
\newcommand{\cataloguescope}[1]{\ifshowcatalogue
 \paragraph{Exact scope in the supplied catalogue.}
 \begin{quote}\small #1\end{quote}\fi}
\hypersetup{pdftitle={Problem 129 - Unconditional classical quantum verification},pdfauthor={Open Problems Math research notebook}}
\fancyhead[L]{\small\sffamily Open Problems Math \textbar\ Research notebook}
\fancyhead[R]{\small\sffamily 129 \textbar\ 27 September 2026}
\numberwithin{equation}{section}
\begin{document}
\setcounter{section}{128}
\setlength{\parskip}{3pt plus 1pt}
\setlist[itemize]{itemsep=3pt,topsep=3pt}
\titlespacing*{\paragraph}{0pt}{1.5ex plus .5ex minus .2ex}{1em}
% ================================================================
% problemId: problem.unconditional-single-prover-classical-verification-of-bqp-computations
% source releaseRank: 129; source releaseStatus: open_with_solved_subcases
% exactTarget SHA-256: ba79ab578a6e8daba3a13423245185a7beeacdb18f4097dd45d4e0cc373905d2
\clearpage
\section{\texorpdfstring{Unconditional Single-Prover Classical Verification of BQP Computations}{Unconditional Single-Prover Classical Verification of BQP Computations}}\label{129}
\subsection{Definitions and mathematical statement}
An interactive proof consists of a verifier and prover exchanging messages, with acceptance determined after polynomially many rounds. Here the verifier is classical probabilistic polynomial time, every message is classical, and the honest prover is uniform quantum polynomial time. For each $L\in\mathsf{BQP}$ require
\[
 x\in L\Rightarrow\Pr[\text{accept with honest prover}]\ge2/3,\qquad
 x\notin L\Rightarrow\sup_{P^*}\Pr[\text{accept with }P^*]\le1/3.
\]
The supremum includes computationally unbounded cheating provers. There is only one prover, and no computational hardness assumption weakens the information-theoretic soundness requirement.
\par\smallskip\textit{Formulation and concept references:} \sref{129}{1}.
\cataloguescope{Determine whether every language L in BQP admits a polynomial-round interactive proof with only classical messages between a probabilistic classical polynomial-time verifier and a single honest uniform quantum polynomial-time prover, with completeness at least 2/3 and soundness at most 1/3 against every computationally unbounded cheating prover.}
\subsection{Short English statement}
Can a fully classical efficient verifier check arbitrary efficient quantum computations using one prover and unconditional security?
% BEGIN RESEARCH ADDITION 129
\subsection{Research attempt}

\paragraph{Current frontier and scope.}
The original issue formulates a single-prover protocol with a classical
efficient verifier, an efficient honest quantum prover, and soundness
against an \emph{unbounded} cheating prover. \sref{129}{1}
Mahadev's classical verification theorem is an important adjacent result,
but its security is computational and uses a quantum hardness assumption;
it is not the information-theoretic statement selected here.
\eref{129}{1}  These distinctions remain part of the target even if
one changes round complexity or amplifies the completeness--soundness
gap.  The literature check, made on 27 September 2026, did not supply
a verified general protocol satisfying all the selected requirements.

\paragraph{Attack map.}
One possible approach is to make the public interaction small enough
that a classical verifier can test all consistency conditions.  The
lemma below shows that a polynomial-size public transcript tree would
actually make the whole language classically decidable with bounded
error.  A second approach is to use local consistency questions to
prevent adaptive cheating; a two-constraint example diagnoses the
need for a genuine consistency mechanism.  A third is to implement a
powerful classical interactive proof using an honest quantum prover;
its missing ingredient is an efficient way for that prover to answer
all auxiliary queries, not just to compute the final decision bit.
None of these obstacles is an unconditional impossibility theorem for
the original polynomial-communication problem.

\paragraph{Attempt I: a small-transcript simulation lemma.}
Fix an input $x$ of length $n$.  Consider an all-classical interaction
with a probabilistic polynomial-time verifier $V$.  Its private random
tape has a polynomial upper length bound and can be sampled in advance.
Suppose the legal public transcript tree has at most $L=\operatorname{poly}(n)$
leaves and can be enumerated in polynomial time.  For example, a fixed
message schedule using at most $C\log_2 n$ communicated bits has this
property.  Message lengths, stopping signals and any visible scheduling
information must be included in this budget; an uncharged timing
channel is not permitted.

For a complete public history $h$, define
\begin{equation}\label{129:leafweight}
 w_h=\Pr_R\left[\begin{gathered}
  V\text{ sends precisely its messages in }h\text{ and accepts,}\\
  \text{when supplied the prover messages in }h
 \end{gathered}\right].
\end{equation}
This is an \emph{unconditional} probability over $R$, not a probability
conditioned on reaching a potentially rare transcript.  It is
classically estimable: sample $R$, simulate $V$ with the designated
prover messages, and output zero if any verifier message or stopping
decision differs from $h$.

For any deterministic prover strategy $\pi$, let $\mathcal L(\pi)$ be
the set of leaves consistent with its choices at prover nodes.  Then
\begin{equation}\label{129:value}
 \Pr[V^\pi(x)\text{ accepts}]=\sum_{h\in\mathcal L(\pi)}w_h.
\end{equation}
The acceptance events on the right are disjoint: a random tape and a
fixed strategy generate only one transcript.  A randomized prover is
a distribution over deterministic contingent strategies and cannot
exceed the largest deterministic acceptance probability.  The same
holds for a prover using quantum internal memory when all messages
are classical.  Its conditional response distributions on classical
histories can be reproduced by an unbounded classical randomized
strategy, and every such behavioral strategy is a mixture of
contingent deterministic choices on this finite tree.

The maximum in \eqref{129:value} is computable from the leaf weights by
a bottom-up recursion.  At a prover node take the \emph{maximum} over
its children; at a verifier node take the \emph{sum}.  No fresh branching
probabilities are multiplied in at verifier nodes: they have already
been included in $w_h$.  Choices after different verifier messages may
be optimized independently, because the prover sees those messages.
Induction on subtrees proves that this recursion equals the optimal
unbounded-prover acceptance probability at the root.

Estimate each $w_h$ to additive error
$\delta=1/(100L)$ with simultaneous success probability at least $0.99$.
Independent Bernoulli sampling and the elementary exponential tail
bound require $O(\delta^{-2}\log(100L))$ samples per leaf, a polynomial
number overall.  For each fixed strategy, replacing $w_h$ by its estimate
changes \eqref{129:value} by at most $L\delta$.  Taking a maximum
preserves this error bound.  Thus the computed optimal acceptance
probability has additive error at most $1/100$, except with probability
at most $0.01$.

\textbf{Small-transcript lemma.}
Any language with such a polynomial-size transcript-tree protocol,
completeness at least $2/3$, and soundness at most $1/3$ against every
prover lies in $\mathsf{BPP}$.

\emph{Proof.} On yes-inputs the optimal acceptance probability is at
least $2/3$, regardless of how the honest strategy is implemented.
On no-inputs the soundness quantifier bounds that optimum by $1/3$.
The estimate above, thresholded at $1/2$, decides the language with
error at most $0.01$ in polynomial time. \hfill$\square$

The proof does not enumerate exponentially many contingent strategies;
the dynamic program optimizes them together.  It does not divide by a
small probability of reaching a history.  These two details are
essential for a valid polynomial-time simulation of private-coin
protocols.

\paragraph{Consequences and the remaining communication gap.}
If every $\mathsf{BQP}$ language had a protocol of the stated kind with
$O(\log n)$ total public communication and a fixed schedule, the lemma
would give $\mathsf{BQP}\subseteq\mathsf{BPP}$; the reverse containment
comes from simulating classical randomness quantumly.  This is a
conditional collapse consequence, not a proved separation that rules
out such protocols.  It flags the depth of any claimed logarithmic
compression of the verification task.

For the actual catalogue problem, polynomially many communicated bits
can give exponentially many leaves.  The same recursion is still a
finite exact characterization, but estimating every leaf to error
$O(1/L)$ is no longer a polynomial-time algorithm.  Replacing $L$ by
its logarithm without a new structural theorem would invalidate both
the error budget and the runtime.  The original problem permits this
large tree; the lemma therefore does not disprove it.

\paragraph{Attempt II: can small local tests enforce one global computation?}
A minimal adversarial example is the inconsistent pair of classical
requirements $b=0$ and $b=1$.  A verifier that asks one uniformly chosen
requirement and lets the prover supply the requested bit accepts an
adaptive prover with probability one.  No fixed bit satisfies both
requirements, but the local response need not have been fixed before
the challenge.  Repeating this interaction without a binding
consistency mechanism does not fix the problem.

This example does not prove an impossibility theorem for classical
interactive proofs.  Such proofs can use algebraic consistency and
other mechanisms much richer than this test.  It does show that the
proposed inference ``every local constraint is checked with positive
probability, hence one global witness exists'' is false.  A claimed
protocol must prove the existence or adequate consistency of a single
underlying computation against an adaptive unbounded strategy.

One can also imagine expressing quantum circuit amplitudes as sums
and asking the prover for intermediate sums.  Exact circuit expansion
has exponentially many paths.  The fact that a quantum computer can
sample the circuit's output distribution does not furnish exact values
of all arbitrary partial path sums requested by a classical verifier.
A complete version of that route needs an efficient honest-response
algorithm and an information-theoretic consistency proof.  Neither
is supplied by writing down the path-sum identity.  Conversely, a
cryptographically binding mechanism does not give the target's
unbounded-prover soundness simply because the honest prover is quantum
polynomial time; the security distinction in \eref{129}{1} is substantive.

\paragraph{Stress test.}
The small-transcript theorem includes verifier private randomness and
adaptive prover responses.  It does not assume that all verifier coins
are public.  Leaf weights can be tiny, since only absolute accuracy is
used.  Acceptance probabilities for different leaves are not incorrectly
normalized as a single strategy-independent distribution: their sum
may exceed one when incompatible prover choices are included.  The
max/sum recursion handles that distinction.

The honest-prover restriction was not needed for the simulation lemma,
so no efficient description of its quantum state is hidden in the
argument.  The classical language formulation is maintained; no
oracle or promise-problem statement is substituted.  The companion
script checks the recursion against explicit enumeration of contingent
strategies on small private-coin transcript trees, including rare
histories, using rational arithmetic.  Those checks do not establish a
polynomial-size transcript tree for arbitrary quantum computations.

\paragraph{Outcome.}
\textbf{Outcome: Exploratory approach with unresolved gap.}

The proved small-transcript simulation places polynomial-size transcript-tree
protocols in $\mathsf{BPP}$.  It does not handle the exponential transcript
trees allowed by the original problem.  A general unconditional protocol
or impossibility theorem remains unproved; no priority claim is made.
% END RESEARCH ADDITION 129
\subsection{Sources}
\begin{itemize}\raggedright
\item[\textnormal{[S1]}]\hypertarget{source:129:1}{} Mario Krenn, Formalize Open Quantum Problem \#49: Single-prover interactive proofs for quantum computations, google-deepmind/formal-conjectures issue 3466, opened 6 March 2026.
\url{https://github.com/google-deepmind/formal-conjectures/issues/3466}.
{\small\textit{Catalogue formulation source.}}
% BEGIN RESEARCH REFERENCES 129
\item[\textnormal{[E1]}]\hypertarget{extra:129:1}{} Urmila Mahadev.
\emph{Classical Verification of Quantum Computations}, arXiv:1804.01082 (2018), abstract and introduction; computational soundness under a hardness assumption.
\url{https://arxiv.org/abs/1804.01082}.
{\small\textit{Consulted 27 September 2026 for the indicated result or scope.}}
% END RESEARCH REFERENCES 129
\end{itemize}


\end{document}
